\chapter{Models}
\label{ch:models}

Four things are learned in this work: the score weights of
Section~\ref{sec:weights}, and the three judges that variants V1, V2 and
V4 use to re-order the score's fifteen best candidates.
Table~\ref{tab:models} summarises them. This chapter describes each, and
gives particular attention to the design decisions that make their
comparison fair, since those decisions are where an earlier version of
this work was weakest.

\begin{table}[H]
\centering
\caption[The four learned components]{The four learned components, their
data and their evaluation protocol. The seven evaluation plans are removed
by a single shared loading function before any dataset is constructed, so
the exclusion is enforced in one place rather than four.}
\label{tab:models}
\small
\begin{tabularx}{\textwidth}{lXll}
\toprule
Component & Data & Split & Repeats \\
\midrule
Score weights & 52,724 paired comparisons over 11,098 plans & by plan & 5 seeds \\
Feature judge (V1) & 24,920 rows: 12,460 real arrangements and their altered copies & by plan & 5 seeds \\
Image judge (V2) & the same rows, rasterised to $64\times64$ & by plan & 3 seeds \\
Quantum judge (V4) & 400 to 19,936 rows, 4 selected measurements & by plan, one fixed test set & 1 per size \\
\bottomrule
\end{tabularx}
\end{table}

\section{Constructing the training examples}
\label{sec:negatives}

A judge learns from examples: real arrangements labelled plausible,
implausible ones labelled otherwise. Since no dataset supplies the second
class, it must be constructed, and the honesty of the judge depends
entirely on how.

\subsection{The circularity that had to be removed}

The previous version of the feature judge constructed its implausible
examples by editing the \emph{measurements}. It multiplied the two
bay-variability figures, overwrote the interior-column proportion,
stretched the aspect ratio, and recalculated the regularity measure
derived from them. Those are four of the six numbers the model then read.
The task therefore reduced to detecting which measurements had been
altered, and a single-threshold rule on one measurement recovered most of
the reported performance. The model was, in effect, grading its own example
generator.

\subsection{Examples constructed in geometry}

Implausible examples are now created by moving the \emph{column positions},
before any measurement or pixel is computed. Four kinds are generated,
cycling evenly, and Figure~\ref{fig:negatives} shows them.

\begin{description}[style=unboxed,leftmargin=0pt]
\item[Transplant.] Another plan's real arrangement, rescaled onto this
plan's footprint. Everything about it is genuinely building-like; it is
simply the wrong building's answer. This tests whether the judge can tell
that an arrangement suits \emph{this} plan.
\item[Jitter.] Every column displaced by a random amount drawn from a
normal distribution with standard deviation 0.35 of the median span. Tests
whether columns align into lines at all.
\item[Shear.] The arrangement skewed so that its lines are no longer
perpendicular. Tests orthogonality.
\item[Dropout.] Thirty per cent of columns removed and some survivors
displaced by a span. Tests the completeness of the load path.
\end{description}

\begin{figure}[H]
\centering
\includegraphics[width=\textwidth]{07_negatives}
\caption[Real and altered arrangements]{One real arrangement and its four
altered counterparts. The alteration is applied to the column
\emph{coordinates}, before any measurement is taken, so the judge must
distinguish real geometry from broken geometry rather than detect which
numbers were edited.}
\label{fig:negatives}
\end{figure}

The same examples and labels are used by all three judges. That shared
basis is what allows the comparison between them to be a comparison of
methods rather than of datasets.

\section{The feature judge}
\label{sec:ml1}

The feature judge is a gradient-boosted decision ensemble, 200 trees of
depth 3, over six measurements chosen to be independent of scale: aspect
ratio, bay variability in each direction, regularity, the proportion of
interior columns, and the ratio of mean bay sizes. The restriction to
scale-independent measurements is forced by the corpus, whose three sources
use three unit systems, as Section~\ref{sec:corpus} explained.

Positives are the derived arrangements of 12,460 plans, with the seven
evaluation plans removed before construction. Splits are grouped by plan,
so that a plan and its altered copy can never fall on opposite sides, and
are repeated over five seeds.

\section{The image judge}
\label{sec:ml2}

The image judge reads a picture of the arrangement. Each is drawn on a
$64\times64$ grid of pixels in two layers: the first holds the
load-bearing walls as single-pixel lines, the second the columns as
three-by-three blocks. Everything drawn is scaled to fit the canvas, which
removes absolute size entirely. Figure~\ref{fig:mlinputs} shows the
inputs. The network is small by design, three convolutional blocks with
pooling, global average pooling, dropout and a single output, totalling
about 15,000 weights, because the corpus contains only some twelve
thousand positive examples.

\begin{figure}[H]
\centering
\includegraphics[width=\textwidth]{11_inputs}
\caption[Inputs to the image judge]{Inputs to the image judge for one
plan. The upper row is the wall layer and the lower row the column layer,
for the real arrangement and its four altered counterparts. The wall layer
is identical across a row; only the registration between columns and walls
changes.}
\label{fig:mlinputs}
\end{figure}

The scaling to a fitted canvas is the design argument for an image judge on
this corpus: a plan in metres and a plan in drawing pixels become the same
kind of image, which sidesteps the unit problem that confines the feature
judge to six measurements. The same property is also its cost. Once
everything is scaled to fit, the judge cannot distinguish a six-metre span
from a six-foot one. This is a deliberate trade rather than an oversight,
and Section~\ref{sec:mlresults} quantifies what it buys.

Implausible examples are created by moving the column positions and
\emph{then} drawing them. Adding noise to the pixels instead would
reproduce, in image form, exactly the circularity that
Section~\ref{sec:negatives} removed from the feature judge.

\subsection{The blanked-wall control}

The comparison between the two judges is confounded by a difference in
information as well as in architecture: the image judge sees the walls and
the feature judge does not. A third model resolves the confound, namely the
same network on the same rows with the wall layer set to zero.
Section~\ref{sec:mlresults} reports all three, and the result changes the
conclusion.

\section{The quantum judge}
\label{sec:quantum}

The quantum judge applies a variational quantum classifier to the same
task. Four measurements are selected by analysis of variance on the
training data and scaled into the interval $[0, \pi]$. Each becomes the
rotation angle of one quantum bit; the bits are then entangled and rotated
by trainable amounts, and one is measured to produce the output, which a
logistic function converts to a probability. Training adjusts the
trainable rotations by gradient descent in the same manner as any neural
network.

Two circuit designs are tested: a standard variational classifier with one
data encoding followed by two strongly entangling layers, and a design that
re-introduces the input data at every layer, which is known to make small
circuits considerably more expressive
\citep{perezsalinas2020reupload}. Both run on a state-vector simulator,
which reproduces an ideal quantum computer with no noise and no
decoherence; real hardware would only be harder.

\subsection{What can and cannot be claimed}

Four simulated quantum bits cannot support a claim of quantum advantage in
either direction, and none is made. Two features of the design make the
comparison informative nonetheless.

\textbf{Capacity-matched controls.} Logistic regression is fitted on the
\emph{same four} measurements, the \emph{same} rows and the \emph{same}
test set as the circuits. Comparing a four-qubit circuit against a large
ensemble trained on everything would judge it unfairly; comparing it
against nothing would tell the reader nothing. Gradient boosting is also
reported on the same four measurements and on all six, to show what
additional capacity and additional information are worth.

\textbf{A training-set size sweep on one fixed test set.} Every model is
fitted at 400, 2,000, 8,000 and 19,936 training rows and scored on the same
held-out plans. Holding the test set fixed is essential rather than
decorative: an interim version of this experiment measured small-sample
models on a small test set and full-corpus models on a larger one, and the
circuits then appeared to \emph{degrade} with more data. The
capacity-matched control fell by the same amount, which is what revealed
the error rather than the finding. That interim result is not reported as a
result.

\subsection{Reaching the full corpus}

An earlier implementation evaluated the circuit once per row in a Python
loop, and it was that, rather than any property of the physics, which
confined this stage to a few hundred examples. Evaluating a whole
minibatch in a single call using parameter broadcasting is approximately
160 times faster and makes the entire corpus reachable on a simulator. The
judge that the pipeline uses is trained on the canonical 80\% training
partition, so that all three judges see identical data and none has seen
the held-out 20\% or the seven demonstration plans.

\section{Combining a judge with the score}
\label{sec:blend}

Each judge receives the fifteen best-scoring candidates and re-orders them
by Equation~\eqref{eq:blend}, repeated here for convenience:
\[
\text{combined} = z(\text{penalty}) + \lambda\bigl(1 - P(\text{plausible})\bigr).
\]
The penalty is standardised across the \emph{whole} candidate set rather
than within the shortlist. This detail matters. Standardising within the
shortlist spreads fifteen near-identical scores across a full interquartile
range, which inflates the score's differences and mutes the judge at any
reasonable $\lambda$; an interim run with that error produced identical
tables for all four variants, which is how it was caught. Standardising
over the full set leaves the shortlist's differences small, so the judge
holds real authority.

The default $\lambda$ is 0.3. Section~\ref{sec:sensitivity} varies it from
0 to 1 and reports the effect, rather than defending the default by
argument.

\section{The four variants}

Table~\ref{tab:variants} states the four paths under comparison. All share
every stage before the branch, so the comparison is controlled: identical
grid, identical candidates, identical score table. Variant V3 is the
baseline. If V1, V2 and V4 do not improve on it under the protocol of
Chapter~\ref{ch:results}, that is the result of the experiment and is
reported as such.

\begin{table}[H]
\centering
\caption[The four pipeline variants]{The four pipeline variants. Each runs
under both sets of score weights, giving the eight rows of
Table~\ref{tab:accuracy}.}
\label{tab:variants}
\small
\begin{tabularx}{\textwidth}{llX}
\toprule
Variant & Judge & Representation the judge reads \\
\midrule
V1 & Feature judge & Six measurements, independent of scale \\
V2 & Image judge & A $64\times64$ two-layer picture, scaled to fit \\
V3 & None & --- (the score's own order stands; the baseline) \\
V4 & Quantum judge & Four measurements as quantum-bit rotation angles \\
\bottomrule
\end{tabularx}
\end{table}
